% Solutions to the problems of Chapter 4 (Learning-Based and Hybrid Sensor Fusion Techniques).
% Numbers reproduced by companion_repo/exercises/ch04/solutions/.
\begin{sol}{prob:ch04:1}
With $g_i\propto\exp(\beta q_i)$ and $q=(0.3,0.5,0.9)$: for $\beta=0$ the gate is uniform, $(1/3,1/3,1/3)$, and the fused estimate is the plain mean $42.50$~m; for $\beta=4$ the gate is $(0.070,\,0.156,\,0.774)$ and the estimate $0.070\cdot41.8+0.156\cdot42.6+0.774\cdot43.1=42.93$~m; for $\beta=20$ the gate is $(0.000,\,0.0003,\,0.9997)$ and the estimate $43.10$~m. Writing $g_i=\exp(\beta(q_i-q_{\max}))/\sum_j\exp(\beta(q_j-q_{\max}))$, every term with $q_j<q_{\max}$ decays as $e^{-\beta(q_{\max}-q_j)}\to0$, so $g\to$ the indicator of the arg-max and $\widehat{x}\to\widehat{x}_{\mathrm{rad}}=43.1$~m. The proportional weights $(0.18,0.29,0.53)$ give $42.7$~m: the softmax with $\beta=4$ is already sharper than the proportional rule, and $\beta$ controls the averaging-to-switching behaviour.
\end{sol}
\begin{sol}{prob:ch04:2}
Reliabilities $p=\sigma(-2+5q)=(0.924,\,0.881,\,0.924)$; information weights $p_i/\sigma_i^{2}=(0.924,\,3.523,\,1.444)$, sum $5.891$; fused estimate $\widehat{x}=(0.924\cdot41.8+3.523\cdot42.6+1.444\cdot43.1)/5.891=42.60$~m with standard deviation $1/\sqrt{5.891}=0.41$~m. Plain weighted least squares gives $42.60$~m and $0.39$~m. With all reliabilities close to one the inflation $\mathbf{R}_i/p_i$ is nearly uniform, so the relative weights, and hence the estimate, coincide with the nominal ones; only the reported standard deviation is slightly larger.
\end{sol}
\begin{sol}{prob:ch04:3}
$\mathrm{ECE}=(100\cdot0.05+200\cdot0.05+300\cdot0.10+250\cdot0.10+150\cdot0.08)/1000=82/1000=0.082$; the detector is over-confident in the three upper bins and under-confident in the lowest one. For $\mathbf{z}=(3,1,0)$: at $T=1$, $\mathbf{p}=(0.844,\,0.114,\,0.042)$; at $T=2$, $\mathbf{z}/T=(1.5,0.5,0)$ and $\mathbf{p}=(0.629,\,0.231,\,0.140)$. The ordering of the logits is preserved by division by a positive scalar, so the arg-max (class 1) is unchanged while the confidence is reduced.
\end{sol}
\begin{sol}{prob:ch04:4}
$l_{\mathrm{occ}}=0.847$, $l_{\mathrm{free}}=-0.405$. Six hits give $L=5.084$ without clamping and $L=3.5$ with clamping. The probability falls below $0.5$ when $L<0$, i.e.\ after $n>5.084/0.405=12.5$, hence $13$ free observations without clamping, and after $n>3.5/0.405=8.6$, hence $9$ free observations with clamping. Clamping bounds the amount of evidence a cell can accumulate, so a vehicle that parks for minutes and then drives away is cleared from the map after a fixed number of free rays instead of a number proportional to how long it was observed; the price is that persistent static structure is held with a bounded confidence.
\end{sol}
\begin{sol}{prob:ch04:5}
The loop-closure weight becomes $\sigma^{-2}=400$. The normal equations read
$\begin{bmatrix}200&-100&0\\-100&200&-100\\0&-100&500\end{bmatrix}\widehat{\mathbf{x}}=\begin{bmatrix}0\\0\\1180\end{bmatrix}$, whose solution is $\widehat{\mathbf{x}}=(0.908,\,1.815,\,2.723)$~m, with odometry residuals $-0.092$~m each and loop-closure residual $+0.023$~m (weighted cost $2.77$). The standard deviation of $\widehat{x}_3$ is $\sqrt{[(\mathbf{A}^{\mathsf T}\mathbf{\Omega}\mathbf{A})^{-1}]_{33}}=0.048$~m. In general, with $n$ odometry steps of variance $\sigma_{\mathrm{odo}}^{2}$ and discrepancy $d$ between dead reckoning and the loop closure, minimising $\sum_{k=1}^{n}\delta_k^{2}/\sigma_{\mathrm{odo}}^{2}+(d-\sum_k\delta_k)^{2}/\sigma_{\mathrm{lc}}^{2}$ over the odometry corrections $\delta_k$ gives by symmetry $\delta_k=\delta$ and $\delta/\sigma_{\mathrm{odo}}^{2}=(d-n\delta)/\sigma_{\mathrm{lc}}^{2}$, i.e.\ $\delta=d\rho/(1+n\rho)$ and loop-closure residual $d-n\delta=d/(1+n\rho)$. For $d=0.3$, $n=3$, $\rho=4$ this is $0.3/13=0.023$~m, as found.
\end{sol}
\begin{sol}{prob:ch04:6}
Running \repopath{companion\_\allowbreak repo/\allowbreak exercises/\allowbreak ch04/\allowbreak solutions/\allowbreak prob6\_\allowbreak miscalibrated\_\allowbreak head.py} (same scenario, a different noise realisation from the example) gives: $c=0.5$: RMSE $0.527$~m, ANIS $7.73$, $60.0\,\%$ of NIS values inside the $95\,\%$ band (over-confident: the filter trusts the detections too much and its gate would reject valid ones); $c=1$: RMSE $0.437$~m, ANIS $1.96$, $94.0\,\%$ inside (calibrated); $c=2$: RMSE $0.375$~m, ANIS $0.50$, $91.0\,\%$ inside (under-confident). The under-confident filter has the lowest RMSE here because the simulated target is exactly constant-velocity and the filter's process noise is conservative, so additional smoothing helps; RMSE therefore cannot reveal the miscalibration, whereas the ANIS, which should be close to $n_y=2$, exposes both the factor-of-four over-confidence ($7.73\approx 2/0.5^{2}$) and the under-confidence ($0.50\approx2/2^{2}$).
\end{sol}
\begin{sol}{prob:ch04:7}
With the seed of \repopath{companion\_\allowbreak repo/\allowbreak exercises/\allowbreak ch04/\allowbreak solutions/\allowbreak prob7\_\allowbreak pose\_\allowbreak graph\_\allowbreak 2d.py}, Gauss--Newton converges in $5$ iterations. The position error of the last pose is $0.246$~m from dead reckoning and $0.106$~m after optimisation; the weighted cost drops from $23.07$ to $1.23$ (four factors of three residuals each, nine unknowns, i.e.\ three redundant scalar measurements, so a cost of about $3$ is expected); the marginal standard deviations of the last pose are $0.185$~m, $0.099$~m, and $0.93^{\circ}$. The loop closure constrains the last pose mainly along the directions it measures precisely, and the correction is propagated backwards to the intermediate poses through the odometry factors.
\end{sol}
\begin{sol}{prob:ch04:8}
Concatenation followed by a linear map gives $\mathbf{f}=\mathbf{W}_1\mathbf{f}_c+\mathbf{W}_2\mathbf{f}_l$. A fixed gate $\mathbf{g}$ is obtained with $\mathbf{W}_1=\diag(\mathbf{g})$, $\mathbf{W}_2=\diag(\mathbf{1}-\mathbf{g})$. An input-dependent gate produces terms such as $\sigma(\mathbf{w}_g^{\mathsf T}[\mathbf{f}_c;\mathbf{f}_l])\,\mathbf{f}_c$, which are nonlinear (at least bilinear) in the inputs and cannot be represented by any fixed linear map, so gating strictly extends concatenation. For the query $(0,1)$ the scaled scores against the keys $(1,0),(0,1),(-1,0)$ are $(0,\,0.707,\,0)$, the softmax weights $(0.248,\,0.503,\,0.248)$, and the output $0.248\cdot2+0.503\cdot0+0.248\cdot(-2)=0$: the query now matches the second key, and the two opposite values cancel. In the example of the chapter the query $(1,0)$ gave weights $(0.576,0.284,0.140)$ and output $0.872$; the attention output depends entirely on the similarity between query and keys.
\end{sol}
